\section{Préschémas réduits~; condition de séparation}
\label{section:I.5-fr}

\subsection{Préschémas réduits.}
\label{subsection:I.5.1-fr}

\begin{proposition}[5.1.1]
\label{I.5.1.1-fr}
Soient $(X,\mathscr{O}_X)$ un préschéma, $\mathscr{B}$ une
$\mathscr{O}_X$-Algèbre quasi-cohérente. Il existe un
$\mathscr{O}_X$-Module quasi-cohérent et un seul $\mathscr{N}$ dont la fibre
$\mathscr{N}_x$ en tout $x\in X$ soit le nilradical de l'anneau
$\mathscr{B}_x$. Lorsque $X$ est affine, et par suite
$\mathscr{B}=\widetilde{B}$, où $B$ est une algèbre sur $A(X)$, on a
$\mathscr{N}=\widetilde{\mathfrak{N}}$, où $\mathfrak{N}$ est le nilradical
de $B$.
\end{proposition}
